\begin{tabular}{>{\raggedright\arraybackslash}p{2.9cm}>{\raggedright\arraybackslash}p{5.2cm}>{\raggedright\arraybackslash}p{2.6cm}>{\raggedright\arraybackslash}p{5.4cm}}
\toprule
问题 & 数据证据 & 处理 & 对数字的影响 \\
\midrule
HistData 时钟 03-29 才切到 UTC−4 & 图 1 两条独立证据 & 分段换算 & 恒 UTC−5：XAU 测试基线 RMSE 29.12（本文 2.28）；按美国夏令时：2.39 \\
XAU 指数 03-05 04:40–05:54 停滞 & 指数−金价最大 30.1bp，永续−金价中位 −3.6bp & 解释为指数异常，不改数据 & 该期官方 −0.1652bp 落在包络 [−1.54, −0.29] 外，是全样本唯一一期 \\
HistData 同分钟多行实为两路报价交错 & 732 分钟、387 行时间倒序；同分钟收盘最大相差 72.6bp & 按题面规则合并 & 影响 46 个样本；改用首行收盘，XAU 测试基线 RMSE 2.275→2.306 \\
日休前最后一分钟金价报价偏低 & 156 次日休：收盘前一分钟基差 7.59→最后一分钟 9.22bp & 保留（规则要求），模型用距收盘特征吸收 & 299 个样本的 t 或 t+4h 落在该分钟 \\
07-16 换算实际在 08:33 生效 & 溢价指数在 08:33 跳 −6.7bp，此前冻结 12 分钟 & 特征只用偏离、不用水平 & 指数腿均值 0.81→6.95bp \\
资金费 06-24 04:00 缺失≠未结算 & 下一期按 4h 窗口误差 0.002bp，按 8h 0.894bp & 剔除该期，不改下一期窗口 & 若误当成未结算、按 8h 窗口算 08:00 期，误差 0.89bp \\
\bottomrule
\end{tabular}
